% =============================================================
%  ONLINE RESOURCE 2 — Analytical Solution Derivation
%  Supporting Information for Hackman & Zhu (2026)
%  "When Inverse Physics-Informed Neural Networks Become
%   Practically Identifiable: A Case Study in Synaptic Dopamine
%   Transport" — Medical & Biological Engineering & Computing
%   submission
% =============================================================
\documentclass[12pt]{article}

\usepackage[margin=1in]{geometry}
\usepackage[utf8]{inputenc}
\usepackage[T1]{fontenc}
\usepackage{lmodern}
\usepackage{amsmath, amssymb, mathtools}
\usepackage{hyperref}
\hypersetup{colorlinks=true, linkcolor=black, citecolor=black,
            urlcolor=blue, pdfborder={0 0 0}}
\usepackage{setspace}
\usepackage{lineno}
\linenumbers
\doublespacing
\setlength{\parindent}{0pt}
\setlength{\parskip}{0.7em}

\title{Online Resource 2: Analytical Solution Derivation \\
\large Supporting Information for ``When Inverse
Physics-Informed Neural Networks Become Practically
Identifiable: A Case Study in Synaptic Dopamine Transport''}

\author{Emmanuel Hackman \and Huiqing Zhu \\
School of Mathematics and Natural Sciences, \\
University of Southern Mississippi, Hattiesburg, MS, USA \\
Corresponding author: \texttt{Emmanuel.Hackman@usm.edu}}

\date{}

\begin{document}
\maketitle

We derive the closed-form solution of the initial-boundary
value problem of equations~(1)--(4) of the main text on the
unconfined planar domain $(x, y) \in \mathbb{R}^2$. This
expression serves as the ground-truth reference for the
forward-problem $L^2$ error reported in Table~2 of the main
text.

The two-dimensional reaction-diffusion equation
\begin{equation}
  \frac{\partial C}{\partial t}
  \;=\; D\, \nabla^2 C \;-\; k\, C,
  \quad \nabla^2 = \partial_{xx} + \partial_{yy},
  \label{eq:S2_pde}
\end{equation}
reduces to a pure diffusion equation under the substitution
\begin{equation}
  C(x, y, t) \;=\; u(x, y, t)\, e^{-k t},
  \label{eq:S2_sub}
\end{equation}
which yields $\partial_t u = D\, \nabla^2 u$ for $u(x, y, t)$
and inherits the same radially symmetric initial profile
$u(x, y, 0) = C_0\, \exp\!\big(-(x^2 + y^2) / (2\sigma^2)\big)$.

The fundamental solution of the diffusion equation on
$\mathbb{R}^2$ is the two-dimensional heat kernel
\begin{equation}
  G(x, y, t) \;=\;
  \frac{1}{4\pi D t}\,
  \exp\!\left( -\frac{x^2 + y^2}{4 D t} \right),
  \label{eq:S2_kernel}
\end{equation}
and the solution with Gaussian initial data follows by
two-dimensional convolution,
\begin{equation}
  u(x, y, t) \;=\;
  \iint_{\mathbb{R}^2}
    G(x - x',\, y - y',\, t)\, u(x', y', 0)\, \mathrm{d}x'\,
    \mathrm{d}y'.
  \label{eq:S2_conv}
\end{equation}
Because both the heat kernel and the radially symmetric
Gaussian initial condition separate as products of independent
$x$ and $y$ factors, the integral evaluates to the product of
two one-dimensional Gaussians,
\begin{equation}
  u(x, y, t) \;=\;
  \frac{C_0\, \sigma^2}{\sigma^2 + 2 D t}\,
  \exp\!\left( -\frac{x^2 + y^2}{2(\sigma^2 + 2 D t)} \right),
  \label{eq:S2_diffusion}
\end{equation}
which describes a radially symmetric Gaussian whose variance
grows linearly in time as
$\sigma_t^2 = \sigma^2 + 2 D t$ in each Cartesian direction,
while the integrated mass
$\iint u\, \mathrm{d}x\, \mathrm{d}y = 2\pi\, C_0\, \sigma^2$
is conserved. Re-applying the substitution~\eqref{eq:S2_sub}
gives the full analytical solution
\begin{equation}
  \boxed{\;
  C(x, y, t) \;=\;
  \frac{C_0\, \sigma^2}{\sigma^2 + 2 D t}\,
  \exp\!\left( -\frac{x^2 + y^2}{2(\sigma^2 + 2 D t)} \right)
  \, e^{-k t}.
  \;}
  \label{eq:S2_full}
\end{equation}

Equation~\eqref{eq:S2_full} is exact on the unconfined domain
$\mathbb{R}^2$. It satisfies the Neumann boundary conditions
$\nabla C \cdot \hat{n} = 0$ on $\partial \Omega$ in the limit
$L \to \infty$, and approximately when $\sigma_t \ll L/2$.
For the parameter values used in this work
($L = 5\,\mu\text{m}$, $\sigma = 0.5\,\mu\text{m}$,
$D = 0.32\,\mu\text{m}^2/\text{ms}$), this condition holds over
the early portion of the simulation window; equation~\eqref{eq:S2_full}
therefore provides a reliable ground-truth reference in that
regime, while finite-domain corrections via the method of images
or a two-dimensional Fourier-cosine expansion become non-negligible
at later times. The finite-difference reference solver of
Section~3 of the main text serves as the complementary baseline
that captures the boundary-driven deviations from~\eqref{eq:S2_full}
at large $t$.

\end{document}
